\subsection{微分形式的模测度}

记 \(\mu\) 为 K 的加法群上的一个 Haar 测度（INT, VII, §1, no. 2）。记 \(\mu^{\otimes n}\) 为测度 \(\mu \otimes \cdots \otimes \mu\) 在 \(K^n\) 上的乘积，并记 \(\mu_{\mathbf{U}}^{\otimes n}\) 为其在 \(K^n\) 的一个开集 U 上的限制。对于 \(a \in K\)，用 mod(a) 表示 a 的模（INT, VII, §1, no. 10 和 AC, VI, §9, no. 1）。若 K = R，则 mod(a) = \textbar a\textbar；若 K = C，则 mod(a) = \textbar a\textbar\^{}2；若 K 是超度量的，则 mod(a) = \(q^{-v(a)}\)，其中 \(q\) 表示 K 的剩余域中元素的个数，\(v\) 是 K 的归一化赋值（AC, VI, §9, no. 1, prop. 1）。

10.1.1. 设 U 和 V 是 \(K^{n}\) 的开子集，且 f: \(U \to V\) 是 \(C^{r}\) 类映射，其中 \(r \in N_{K}\)。设 \(x \in U\)；称 f 在 x 处的雅可比为 \(\operatorname{Jac}_{x}(f)\)，即 f 在 x 处的导数这一线性映射 \(Df(x)\) 的行列式。记 \(\operatorname{Jac}(f)\) 为函数 \(x \mapsto \operatorname{Jac}_{x}(f)\)；函数 \(\operatorname{mod}(\operatorname{Jac}(f))\) 是 U 上取正实值的连续函数。

10.1.2（“积分中的变量替换”）。在 10.1.1 的假设下，设 \(f\) 是从流形 U 到流形 V 的 \(\mathbf{C}^{r}\) 类同构。则测度 mod(Jac(f)). \(f\) 在 \(\mu_{\mathrm{U}}^{\otimes n}\) 下的像是 \(\mu_{\mathrm{V}}^{\otimes n}\)；对于每个具有紧支集的连续函数 \(\varphi: \mathrm{V} \to \mathbf{C}\)，有

\[
\int_ {\mathrm{V}} \varphi (y) \mu^ {\otimes n} (y) = \int_ {\mathrm{U}} \varphi (f (x)). \mathrm{mod} (\operatorname{Jac} _ {x} (f)) \mu^ {\otimes n} (x).\tag{1}
\]

10.1.3. 设 X 是一个 \(\mathbf{C}^r\) 类流形，A 是 X 的子集。若对于 X 的每个图 \((\mathrm{U},\varphi ,\mathrm{K}^n)\)，集合 \(\varphi (\mathrm{A}\cap \mathrm{U})\) 是 \(\mu^{\otimes n}\)-零测的（INT, IV, § 2, no. 2），则称 A 局部零测；这一条件不依赖于 \(\mu\) 的选择，只需对其定义域覆盖 A 的一族图验证即可。包含在可数个局部零测集并集中的每个集合都是局部零测的。

例子

\begin{enumerate}
\def\labelenumi{\alph{enumi})}
\item
  X 中每个在每一点都具有余维 \(\geqslant1\)（即内点为空的）子流形，都是局部零测的。
\item
  设 g 是 X 上的 \(C^{r}\) 类函数，且 \(A_{g}\) 是满足 \(x \in X\) 的点 \(g(x) = 0\) 的集合；假设 \(A_{g}\) 的内点为空，且 \(r = \omega\)；则 \(A_{g}\) 局部零测。
\item
  设 \(f: Y \to X\) 是 \(C^{r}\) 类流形的态射，其中流形 \(Y\) 是紧集的可数并。假设不等式 \(\dim_{y} Y \leqslant \dim_{f(y)} X\) 对所有 \(y \in Y\) 成立。在 \(f\) 下，\(Y\) 的每个局部零测子集的像都是 \(X\) 的局部零测子集。在 \(f\) 下，\(f\) 不平展的点集的像也具有同样性质（“第一 Sard 定理”）；特别地，若不等式 \(\dim_{y} Y < \dim_{f(y)} X\) 对所有 \(y \in Y\) 成立，则 \(f(Y)\) 是 \(X\) 的局部零测子集。
\item
  假设 K 的特征为零。设 \(f\colon Y\to X\) 是 \(\mathbf{C}^{r}\) 类流形的态射（\(r\in N_{K}, r\geqslant\infty\)），其中 Y 是紧集的可数并。令 C 为 Y 中 \(f\) 不是淹没的点的集合（“临界点”）。则 \(f(\mathbf{C})\) 是 X 的局部零测子集（“第二 Sard 定理”）。\(^{(1)}\)
\end{enumerate}

10.1.4. 设 X 是 C' 类仿紧流形。在 X 上存在且仅存在一个等价测度类 M（INT, V, § 5, n° 6），使得对于每个 \(\nu \in \mathbf{M}\) 以及 X 的每个图 \(c = (\mathrm{U}, \varphi, \mathrm{K}^n)\)，\(\varphi\) 在 U 上的限制经 \(\nu\) 映射后的像等价于 \(\mu_{\varphi(\mathrm{U})}^{\otimes n}\)。称 M 为 X 上的典范测度类。对于 X 的子集 A，要使其局部零测（10.1.3），充要条件是：对于某一个（分别每一个）测度 \(\nu \in \mathbf{M}\)，A 是局部零测的（INT, IV, § 2, n° 2）。

若 K = R 且 \(r \neq \omega\)（分别 K = R 或 C 且 \(r = \omega\)，分别 K 是超度量的），则存在 \(\nu \in M\)，使得对于 X 的每个图 \(c = (\mathrm{U}, \varphi, \mathrm{K}^{n})\)，\(\varphi(\nu_{\mathrm{U}})\) 关于 \(\mu_{\varphi(\mathrm{U})}^{\otimes n}\) 的密度处处 \textgreater0 且为 \(C^{r-1}\) 类（分别为 \(C^{\infty}\) 类，分别局部常数）。若 \(\nu\) 和 \(\nu'\) 是两个这样的测度，则 \(\nu'\) 关于 \(\nu\) 的密度处处 \textgreater0 且为 \(C^{r-1}\) 类（分别为 \(C^{\infty}\) 类，分别局部常数）。

10.1.5. 设 X 是 \(\mathbf{C}^r\) 类流形；令 \(\mathrm{T}'(\mathrm{X})\) 为 X 的余切丛（8.2.2），并置 \(\Omega = \det(\mathrm{T}'(\mathrm{X}))\)（7.9.9）。向量丛 \(\Omega\) 在每一点的秩为 1；当 X 是纯 \(n\) 维时，它与向量丛 \(\operatorname{Alt}^n(\mathrm{T}(\mathrm{X}); \mathrm{K}_{\mathbf{x}})\) 等同。令 \(\omega\) 为 X 上 \(\Omega\) 的一个截面；虽属术语上的滥用，称 \(\omega\) 为 X 上的最高次数微分形式。设 \(c = (\mathrm{U}, \varphi, \mathrm{K}^n)\) 是 X 的一个图；记 \(u^1, \ldots, u^n\) 为 \(\mathrm{K}^n\) 上的坐标函数。在 \(f_c\) 上存在且仅存在一个函数 \(\varphi(\mathrm{U})\)，使得 \(\omega|_{\mathrm{U}} = \varphi^*(f_c.du^1 \wedge \cdots \wedge du^n)\)。若对于 X 的每个图 \(\omega\)，实值函数 mod \(c = (\mathrm{U}, \varphi, \mathrm{K}^n)\) : \((f_c)\) 关于 \(\varphi(\mathrm{U}) \to \mathbf{R}\) 局部可积（INT, IV, § 4, n° 1），则称 \(\mu_{\varphi(\mathrm{U})}^{\otimes n}\) 具有局部可积模；只需对 X 的一个图册中的图验证这一性质。若 \(\omega\) 连续，特别地，若 \(\omega\) 是 \(\mathbf{C}^s\) 类的，其中 \(s \in \mathbb{N}_{\mathbf{K}}, s \leqslant r - 1\)，则 \(\omega\) 具有局部可积模。

10.1.6（“由最高次数微分形式定义的正测度”）。保留 10.1.5 的记号，并进一步假设 X 是分离的，且 \(\omega\) 具有局部可积模。若 \(c = (\mathrm{U}, \varphi, \mathrm{K}^n)\) 是 X 的一个图，记 \(\nu_c\) 为 \(\varphi(\mathrm{U})\) 上由测度 \(\mu_{\varphi(\mathrm{U})}^{\otimes n}\) 与函数 mod( \(f_c\) ) 的乘积得到的测度（INT, V, § 5, n° 2, def. 2），并令 \(\alpha_c\) 为 \(\nu_c\) 经 \(\varphi^{-1}\) 映射后的像。在 X 上存在且仅存在一个测度 \(\alpha\)，使得对于 X 的每个图 \(c = (\mathrm{U}, \varphi, \mathrm{K}^n)\)，\(\alpha\) 在 U 上的限制等于 \(\alpha_c\)。称测度 \(\alpha\) 为 \(\omega\) 的模，记为 mod( \(\omega\) ) \(_\mu\)。它是一个正测度。当 K = R 且 \(\mu\) 是 Lebesgue 测度时，用 \(|\omega|\) 代替 mod( \(\omega\) ) \(_\mu\)。

若 X 是纯 n 维的，且 a 是实数 \textgreater0，则有 \(\text{mod}(\omega)_{a\mu} = a^{n} \text{ mod}(\omega)_{\mu}\)。

设 A 是 X 的子集。要使 A 局部零测（10.1.3），充要条件是：对于 X 的每个开集 U 以及 U 上每个具有局部可积模的最高次数微分形式 \(\omega\)，集合 A ∩ U 关于测度 mod( \(\omega\) ) \(_{\mu}\) 局部零测。

此外，假设 X 是仿紧的，并设 \(\nu\) 是属于典范等价类 M（10.1.4）的 X 上测度。测度 mod \((\omega)_{\mu}\) 基于 \(\nu\)（INT, V, § 5, no. 2, def. 2）。要使 mod \((\omega)_{\mu}\) 属于 M，充要条件是满足 \(x \in \mathbf{X}\) 的 \(\omega(x) = 0\) 的集合是局部零测的。

10.1.7. 例。--- 设 G 是 K 上的有限维、维数为 n 的 Lie 群，且 \(\omega\) 是 G 上的 n 次微分形式，在左平移下不变且非零。相应的测度 \(\text{mod}(\omega)_{\mu}\) 是局部紧群 G 上的左 Haar 测度。

当 G 是乘法群 \(K^{*}\) 时，可取 \(\omega\) 为形式 dx/x，并有 \(\text{mod}(\omega)_{\mu} = (\text{mod})^{-1}.\mu\)，cf.~INT, VII, § 1, no. 10, prop. 14。
\begin{enumerate}
\def\labelenumi{(\arabic{enumi})}
\setcounter{enumi}{6}
\tightlist
\item
\end{enumerate}
% semantic-fragments: legacy-input-seam
\relax{}
